% ===================================================================
%  नकसा नं. ३ — बचपन आ जवानी।
%
%  Traduction, faite en 2026, en bhojpouri d'aujourd'hui, sur l'ido
%  de texto/io. Regles de la colonne : voir 10-nakasa-01.tex. DEUX
%  scenes, deux numerotations qui repartent de 1 (cles t03-c1-* et
%  t03-c2-*), comme dans le fac-simile. 54 blocs, 129 renvois.
%
%  TROIS SALUTS REFUSES POUR UN SEUL « BON JORNO » (68). Le petit
%  garcon poli salue les gens qu'il connait, et le bhojpouri a trois
%  formules toutes pretes. « राम-राम » est le salut de la plaine, on
%  l'entend d'un bout a l'autre du Bhojpour — mais il nomme un dieu,
%  et la scene est un bapteme catholique dans un village francais.
%  « परनाम » a le meme defaut, plus discret. « सुप्रभात » n'a aucune
%  religion, mais c'est du hindi d'ecole, exactement le registre
%  contre lequel cette colonne est ecrite. La colonne fabrique donc
%  le calque avec ses propres mots — « नीक भोर » — et note ici que ce
%  n'est pas de la timidite : ON MODERNISE LA LANGUE, JAMAIS LES
%  CHOSES, et un salut porte une chose autant qu'un mot.
%
%  LES MOIS SONT ANGLAIS, LES JOURS ONT TROIS MILLE ANS, ET LA
%  COLONNE GARDE LES DEUX. मार्च, अप्रैल, मई — les mois du livret
%  sont ceux du calendrier gregorien, arrives en bhojpouri par
%  l'anglais. La region compte pourtant ses propres mois, चइत et
%  बइसाख, et « le printemps, mars, avril ou mai » se dirait tout
%  naturellement चइत-बइसाख a Chapra. On ne l'ecrit pas : la planche
%  montre un village francais, et un village francais n'a pas de
%  चइत. En revanche सोमार, मंगर, बुध — les jours — sont les vieux
%  noms planetaires, et ceux-la ne changent pas d'un continent a
%  l'autre. Meme alinea, deux couches, et la regle tranche la meme
%  chose dans les deux cas.
%
%  C'EST LA MEME DECISION QUI COMMANDE LA CORRECTION DE « julio ».
%  Le fac-simile ido oublie le gras sur « julio » au bloc
%  t03-c2-01-1, seul mois nu au milieu de deux mois gras.
%  gravuri/korekti.json portait deja le triple bhojpouri AVANT que ce
%  fichier existe — il a ete pose au cablage, pas ici. Le compte doit
%  donc bouger : html.py annoncait « corrections declarees : 63 »
%  hier, et doit annoncer 64 des ce tableau. La colonne n'a pas eu a
%  s'apercevoir de rien.
%
%  « पादरी » LE CURE (15) EST PORTUGAIS, ET C'EST LE MEME MOT QUE LA
%  COLONNE GUJARATIE. Elle avait releve પાદરી parmi ses onze mots
%  portugais, arrives par Diu et Daman. Celui-ci est a mille cinq
%  cents kilometres de la cote, dans la plaine du Gange, et personne
%  n'y tient le mot pour etranger. Avec मेज la table (tableau 1), le
%  releve a maintenant deux preuves que la couche portugaise n'est
%  pas restee sur le littoral.
%
%  LIRE parigi.py NE SUFFIT PAS : IL FAUT ECRIRE LE BLOC EN LE
%  TENANT. Le tableau 1 comptait six paires a retourner et il est
%  sorti a zero divergence du premier jet. Celui-ci en compte
%  vingt-deux, reparties sur dix-huit blocs, et le premier jet en a
%  manque SIX — t03-c1-02-1, -04-1, -09-1, -10-1, -14-1 et
%  t03-c2-10-1. L'outil avait ete lu, la liste etait sous les yeux,
%  et cela n'a rien empeche : au-dela de quelques paires, la lecture
%  prealable devient un survol, et c'est la main qui repose le
%  genitif bhojpouri devant son regi sans y penser. Les six fautes
%  ont toutes exactement la meme forme — le nuage devant le ciel, la
%  fleche derriere le clocher, la ruche devant l'abeille, le mur
%  devant le pot, le quai devant la joute. LE TABLEAU 1 N'AVAIT DONC
%  RIEN PROUVE : il avait six paires, pas vingt-deux.
%
%  ET LA COLONNE RECOPIE UNE BEVUE D'IMPRIMEUR SANS POUVOIR RIEN
%  DIRE. Au bloc t03-c1-09-1 l'ido appelle deux fois le renvoi (18)
%  — les arbres fruitiers, puis les pechers — et donne aux
%  abricotiers le (19), qui est deja le presbytere quatre blocs plus
%  haut. La suite est donc 20, 18, 18, 19. Le controle compare la
%  traduction a l'ido et non a la planche : il exige la bevue et la
%  valide. UN CONTROLE QUI COMPARE DEUX COLONNES NE PEUT PAS SAVOIR
%  QU'ELLES ONT TORT ENSEMBLE.
% ===================================================================

%%K t03-apar-1 apar
\VUsaut{3.0mm}
\VUtitre{15.0pt}{60}{नकसा नं. ३}
\VUsaut{1.6mm}
\VUtitre{11.4pt}{0}{बचपन आ जवानी।}
\VUsaut{3.4mm}

%%K t03-apar-2 apar
(तीसरका आ चउथका नकसा एह तरह से जोड़ल गइल बा कि ऊ चार गो
\VUgras{घर-परिवार के दृस्य} में \VUgras{मौसम}, \VUgras{उमिर},
\VUgras{परिवार}, \VUgras{कपड़ा} आ \VUgras{हैसियत} के जरूरी बात
देखावेला।)

\VUsaut{4.0mm}
%%K t03-tit-1 sub
\VUcentre{12.0pt}{0}{\textit{पहिलका दृस्य।}}
\VUsaut{3.0mm}

%%K t03-tit-2 sub
\VUpk{10.2pt}{40}{गाँव में बपतिसमा।}
\VUsaut{2.4mm}

%%K t03-tit-3 sub
\VUpk{10.2pt}{40}{बसंत।}
\VUsaut{3.0mm}

%%K t03-c1-01-1 p
१. --- हमनी के \VUgras{बसंत} में बानी जा, \VUgras{मार्च},
\VUgras{अप्रैल} भा \VUgras{मई} महीना में, \VUgras{हफ्ता} के दिन में,
\VUgras{सोमार}, \VUgras{मंगर} भा \VUgras{बुध} के। भोर के \VUgras{दस
बजल} बा।

%%K t03-c1-02-1 p
२. --- \VUgras{मौसम} सुंदर बा, \VUgras{हवा} साफ बा। सूरज चमकत बा। कुछ
\VUgras{बदरा} \textsuperscript{(2)} नीला \VUgras{आसमान}
\textsuperscript{(1)} में उठत बा।

%%K t03-c1-03-1 p
३. --- \VUgras{गाछ} पर मुसकिल से \VUgras{पत्ता} बा। पातर
\VUgras{सफेदा} \textsuperscript{(3)} आ ऊँच \VUgras{चिनार} \textsuperscript{(17)}
पर अबहीं ले खाली कोंपल बा।

%%K t03-c1-04-1 p
४. --- \VUgras{अबाबील} \textsuperscript{(4)} लवट आइल बाड़ी सँ आ खुसी से
चहचहात ओह \VUgras{सिखर} \textsuperscript{(7)} के चारो ओर उड़त बाड़ी सँ
जवन \VUgras{घंटाघर} \textsuperscript{(5)} के ह, आ ऊ घंटाघर गाँव के
\VUgras{गिरजाघर} \textsuperscript{(6)} के माथे पर ठाढ़ बा।

%%K t03-tit-4 sub
\VUsaut{3.4mm}
\VUpk{10.2pt}{40}{गिरजाघर।}
\VUsaut{3.0mm}

%%K t03-c1-05-1 p
५. --- ई गिरजाघर बहुते सादा बा, आपन \VUgras{बड़का फाटक}
\textsuperscript{(13)} आ मोट \VUgras{घड़ी} \textsuperscript{(8)} के साथे।
\VUgras{छोटकी सूई} \textsuperscript{(9)} X के अंक पर बा, ऊ \VUgras{घंटा}
देखावेले। \VUgras{बड़की} \textsuperscript{(10)} XII पर बा (दुपहरिया) ; ऊ
\VUgras{मिनट} देखावेले।

%%K t03-c1-06-1 p
६. --- घंटाघर में \VUgras{घंट} \textsuperscript{(11)} खुसी से बाजत बा।
छत के चोटी पर पाथर के छोट \VUgras{क्रूस} \textsuperscript{(12)} ठाढ़ बा।

%%K t03-c1-07-1 p
७. --- गिरजाघर में खाली \VUgras{बड़का फाटक} \textsuperscript{(13)} ना
बा, ओहमें बगल में एगो छोट केवाड़ भी बा। एही छोट केवाड़ से
\VUgras{पादरी} \textsuperscript{(15)} साहेब, आपन \VUgras{चोगा} पहिनले,
\VUgras{पूजा-कोठरी} \textsuperscript{(14)} से निकल के \VUgras{पादरी के
घर} \textsuperscript{(19)} ओर जात बाड़ें।

%%K t03-c1-08-1 p
८. --- ओकरा लगे, इहाँ \VUgras{दूधवाली} \textsuperscript{(24)} आपन
\VUgras{गदहा} \textsuperscript{(23)} के साथे बाड़ी। ऊ सहर से लवटत बाड़ी,
जहाँ ऊ बच्चा लोग खातिर नीक दूध ले गइल रहली।

%%K t03-tit-5 sub
\VUsaut{4.0mm}
\VUpk{10.2pt}{40}{फर-बगइचा।}
\VUsaut{3.4mm}

%%K t03-c1-09-1 p
९. --- अलिबोरों बाबू (गदहा) एह भोर के दउड़ से पूरा खुस बाड़ें। ऊ मजा
लेके ओह हवा में साँस लेत बाड़ें जवना के ओह \VUgras{फूल}
\textsuperscript{(20)} के महक बसवले बा जवन बगल के \VUgras{फर-बगइचा} के
\VUgras{फर के गाछ} \textsuperscript{(18)} के ढँकले बा। ऊ \VUgras{आड़ू के गाछ}
\textsuperscript{(18)} आ \VUgras{खुबानी के गाछ} \textsuperscript{(19)} पूरा
गुलाबी आ उजर बाड़ें, आ ओह लोग से नीक फसल के आस बा।

%%K t03-c1-10-1 p
१०. --- ओही बेरा छोट \VUgras{मधुमक्खी} \textsuperscript{(22)}, जे
\VUgras{छत्ता} \textsuperscript{(21)} के ह, भनभनात ओहिजा आपन मीठ
\VUgras{मधु} बटोरे जात बाड़ी सँ।

%%K t03-c1-11-1 p
११. --- बाकिर ई का होत बा ? इहाँ गाँव के लइका दउड़ के आवत बाड़ें आ डेराइल
\VUgras{हंस} \textsuperscript{(25)} के भगावत बाड़ें। ई त नोटरी के बेटा के
\VUgras{बपतिसमा} के बाद गिरजाघर से जुलूस के निकलल ह।

%%K t03-c1-12-1 p
१२. --- छोटका याकोबुस आपन \VUgras{पलना} \textsuperscript{(26)} में
घंट के खुसी भरल आवाज से जाग जात बा, आ ओकर बहिन मगदलेना, जे खुदे भी
बपतिसमा के जुलूस देखल चाहत बाड़ी, आपन माई से बिनती करत बाड़ी कि घर के
ओसारा पर ओकर बाल झार दीहें आ कपड़ा पहिना दीहें।

%%K t03-c1-13-1 p
१३. --- माई मान जाली। ऊ आपन \VUgras{ओढ़नी} \textsuperscript{(28)}, आपन
\VUgras{चोली} \textsuperscript{(29)} आ आपन \VUgras{एप्रन} \textsuperscript{(30)}
पहिनेली, आ केवाड़ के सामने एगो छोट पीढ़ा पर बइठ जाली। मगदलेना लगे खाली
आपन \VUgras{घघरी} \textsuperscript{(31)} आ आपन \VUgras{कँचुकी}
\textsuperscript{(32)} बा।

%%K t03-c1-14-1 p
१४. --- ऊ कतना खुस बाड़ी ! ऊ आपन मनपसंद फूल भुला जात बाड़ी — आपन
\VUgras{कार्नेशन} \textsuperscript{(34)}, जवना पर \VUgras{तितली}
\textsuperscript{(35)} बइठत बाड़ी सँ, आपन \VUgras{ट्यूलिप}
\textsuperscript{(36)}, आपन \VUgras{मई के फूल} \textsuperscript{(37)},
जवन \VUgras{गमला} \textsuperscript{(38)} में, \VUgras{भीत}
\textsuperscript{(33)} पर बा, आ आपन \VUgras{गुलाब} \textsuperscript{(39)}
जवन एतना सुंदर बाड़ बनवले बा। ऊ जुलूस के मेहरारू लोग के महँग कपड़ा
निहारत बाड़ी।

%%K t03-c1-15-1 p
१५. --- गाँव के लइकन के कपड़ा से जादे मतलब नइखे। ऊ लोग
\VUgras{मिठाई} \textsuperscript{(55)} भा \VUgras{सिक्का}
\textsuperscript{(52)} पावे खातिर दउड़त आ एक-दोसरा के ठेलत बाड़ें। ओहमें
से एगो रोवत बा \textsuperscript{(40)}।

%%K t03-c1-16-1 p
१६. --- दोसरका \VUgras{कुकुर} \textsuperscript{(42)} के गोड़ पर पाँव धर
देलस, आ ऊ चिचिआत भागल आ \VUgras{मुरगी} \textsuperscript{(43)} आ ओकर
\VUgras{चूजा} \textsuperscript{(44)} के भगा देलस।

%%K t03-tit-6 sub
\VUsaut{5.0mm}
\VUpk{10.2pt}{40}{बपतिसमा के जुलूस।}
\VUsaut{4.0mm}

%%K t03-c1-17-1 p
१७. --- जुलूस के आगे \VUgras{धाय} \textsuperscript{(45)} चलत बाड़ी।
ओकरा लगे लमहर फीता वाला सुंदर \VUgras{टोपी} \textsuperscript{(46)} बा। ऊ
\VUgras{नवजात} \textsuperscript{(47)} के गोद में लेले बाड़ी, जे पूरा
\VUgras{झालर} \textsuperscript{(48)} में लपेटल बा।

%%K t03-c1-18-1 p
१८. --- धाय के पाछे \VUgras{धरम-पिता} \textsuperscript{(49)} आ
\VUgras{धरम-माई} \textsuperscript{(53)} आवत बाड़ें। ऊ लोग मिठाई आ सिक्का
बाँटत बाड़ें। धरम-पिता लगे \VUgras{दस्ताना} \textsuperscript{(51)} आ
\VUgras{ऊँच टोपी} \textsuperscript{(50)} बा। धरम-माई हाथ में सुंदर
\VUgras{फूल के गुच्छा} \textsuperscript{(54)} पकड़ले बाड़ी।

%%K t03-c1-19-1 p
१९. --- ओह लोग के पाछे हमनी के बच्चा के \VUgras{चाचा}
\textsuperscript{(56)} आ \VUgras{चाची} \textsuperscript{(58)} के देखत बानी
जा। चाची लगे सुंदर \VUgras{हैट} \textsuperscript{(59)} बा, चाचा लगे करिया
\VUgras{लमहर कोट} \textsuperscript{(57)} आ उजर वास्कट बा। ओकरा बाद
\VUgras{चचेरा भाई} \textsuperscript{(60)} आ \VUgras{चचेरी बहिन}
\textsuperscript{(61)} बाड़ें। ई चचेरी बहिन हाथ में नवजात के
\VUgras{बहिन} \textsuperscript{(62)}, छोटकी बेर्ता के पकड़ले बाड़ी।

%%K t03-c1-20-1 p
२०. --- बेर्ता के दोसरका \VUgras{भाई} \textsuperscript{(63)} पाछे चलत
बा। ऊ एगो \VUgras{रिस्तेदारिन} \textsuperscript{(64)} के हाथ पकड़ले बा।
ओकरा बाद परिवार के \VUgras{संगी} \textsuperscript{(65)} लोग आवत बाड़ें।

%%K t03-tit-7 sub
\VUsaut{4.6mm}
\VUpk{10.2pt}{40}{देखे वाला लोग।}
\VUsaut{3.4mm}

%%K t03-c1-21-1 p
२१. --- जुलूस लगे, इहाँ एगो \VUgras{भिखारी} \textsuperscript{(66)} आपन
\VUgras{बैसाखी} \textsuperscript{(67)} पर टेकल बा। ऊ भीख माँगत बा।

%%K t03-c1-22-1 p
२२. --- दोसरा ओर एगो बहुते \VUgras{सभ्य} \textsuperscript{(68)} छोटका
लइका बा। ऊ जेकरा के चिन्हेला ओकरा के परनाम करेला आ « \VUgras{नीक भोर}
» कहेला।

%%K t03-c1-23-1 p
२३. --- गाँव के लोग बपतिसमा के जुलूस देखे खातिर आपन घर से निकलल बा।
हमनी के एगो \VUgras{किसान} \textsuperscript{(69)} के आपन \VUgras{सूती
टोपी} के साथे आ ओकरा बगल में \VUgras{किसानिन} \textsuperscript{(70)} के
देखत बानी जा। ओकरा बाद \VUgras{बुढ़िया} \textsuperscript{(71)} आपन
\VUgras{लाठी} \textsuperscript{(72)} पर टेकल बाड़ी। आ आखिर में
\VUgras{पनचक्की वाला} \textsuperscript{(73)} आपन बड़का \VUgras{फेल्ट के
टोपी} \textsuperscript{(74)} के साथे, आ एगो जवान किसानिन जे
\VUgras{लकड़ी के जूता} \textsuperscript{(75)} पहिनले आपन छोट बच्चा के
चलल सिखावत बाड़ी।

%%K t03-c1-24-1 p
२४. --- ई दृस्य बहुते मजेदार बा।

\VUsaut{4.0mm}
%%K t03-tit-8 sub
\VUcentre{12.0pt}{0}{\textit{दुसरका दृस्य।}}
\VUsaut{3.0mm}

%%K t03-tit-9 sub
\VUpk{10.2pt}{40}{सरबजनिक परब।}
\VUsaut{2.4mm}

%%K t03-tit-10 sub
\VUpk{10.2pt}{40}{नाव के खेल आ दउड़।}
\VUsaut{3.0mm}

%%K t03-c2-01-1 p
१. --- \VUgras{गरमी} आ गइल। \VUgras{जून} महीना के (जुलाई भा
\VUgras{अगस्त}) चिलचिलात घाम जवान लोग के नहाए आ नाव खेवे खातिर उकसावत
बा।

%%K t03-c2-02-1 p
२. --- नदी के दाहिने किनारे, नाव खेवे वाला लोग नाव के खेल आ दउड़ में
हिस्सा लेवे खातिर कपड़ा उतारत बाड़ें।

%%K t03-c2-03-1 p
३. --- ओहमें से एगो आपन \VUgras{जूता} \textsuperscript{(1)} उतारत बा ;
ऊ ओकर बंद खोलत बा। ओकर दू गो संगी पहिलही से तइयार बाड़ें। अबहीं ओह लोग
लगे खाली आपन \VUgras{बुनल गंजी} \textsuperscript{(2)} आ \VUgras{नहाए के
जाँघिया} \textsuperscript{(3)} बा। ऊ लोग आपन दोस्तन के अगोरत बतियावत
बाड़ें। ऊ लोग छुट्टी के आ ओह परब के बात करत बाड़ें जवन दोसरा किनारे
होत बा।

%%K t03-c2-04-1 p
४. --- ओह लोग लगे भुइयाँ पर ओह लोग के संगी लोग के \VUgras{कपड़ा} धइल
बा : एक जोड़ी \VUgras{बूट} \textsuperscript{(4)}, \VUgras{जूता}
\textsuperscript{(5)}, \VUgras{मोजा} \textsuperscript{(6)},
\VUgras{फेल्ट} \textsuperscript{(7)} आ \VUgras{पुआल}
\textsuperscript{(8)} के टोपी आ \VUgras{टाई} \textsuperscript{(9)}।

%%K t03-c2-05-1 p
५. --- ओह जवान लोग में से एगो आपन \VUgras{कपड़ा} \textsuperscript{(10)}
सलीका से तह कइले बा। ऊ ओकरा के एगो अँगोछा पर धरत बा, जवना में ओकर
पड़ोसी थोड़े देर में दुनो कपड़ा बान्ह ली।

%%K t03-c2-06-1 p
६. --- छठवाँ लइका देर करत बा। ओकरा कपार पर अबहीं ले आपन \VUgras{कैप}
\textsuperscript{(11)} बा। ऊ ना आपन \VUgras{कमीज} \textsuperscript{(12)}
उतरले बा, ना आपन \VUgras{कॉलर} \textsuperscript{(13)}, ना आपन
\VUgras{कफ} \textsuperscript{(14)}। ऊ हाथ में आपन \VUgras{वास्कट}
\textsuperscript{(17)} पकड़ले बा आ ओकरा लगे अबहीं ले आपन
\VUgras{पतलून} \textsuperscript{(16)} आ \VUgras{गैलिस} \textsuperscript{(15)}
बा। ऊ अगिला दउड़ खातिर पक्का तइयार ना हो पाई।

%%K t03-c2-07-1 p
७. --- ओह जवान लोग लगे एगो पगडंडी, जेकरा किनारे बहुते \VUgras{कँटीला
घास} \textsuperscript{(18)} बा, नदी के तीर तक जाले, जहाँ याकोबुस बाबू
\VUgras{नरकट} \textsuperscript{(19)} के बीच ओह \VUgras{आतिसबाजी}
\textsuperscript{(20)} के तइयारी करत बाड़ें जवन साँझ के छोड़ल जाई।

%%K t03-tit-11 sub
\VUsaut{4.4mm}
\VUpk{10.2pt}{40}{दउड़।}
\VUsaut{3.4mm}

%%K t03-c2-08-1 p
८. --- नदी पर कुछ मेहरारू लोग, आपन छाता से आपना के बचावत,
\VUgras{नाव} \textsuperscript{(21)} पर घूमत बाड़ी सँ। ऊ लोग ओह दउड़ के
देखत बाड़ी सँ जवन बहुते मजेदार बा। हर \VUgras{डोंगी} \textsuperscript{(22)}
में चार गो \VUgras{खेवइया} \textsuperscript{(24)} आपन पूरा ताकत से खेवत
बाड़ें, आ ओही बेरा \VUgras{मल्लाह} \textsuperscript{(26)} \VUgras{पतवार}
\textsuperscript{(25)} थामले नाजुक डोंगी के चलावत बा। \VUgras{डाँड़}
\textsuperscript{(23)} ताल से पानी पीटत बा आ हलुक नाव चक्कर आवे वाला
तेजी से चलत बाड़ी सँ।

%%K t03-c2-09-1 p
९. --- कुछ जवान लोग, पुल के खंभा लगे, \VUgras{फिसलन वाला मस्तूल}
\textsuperscript{(28)} के इनाम खातिर होड़ लगावत बाड़ें। ओहमें से एगो
लचकत आ फिसलत मस्तूल पर सँभल के चलत बा कि छोर पर बान्हल छोट
\VUgras{झंडी} \textsuperscript{(29)} तक पहुँच जाव। ओकर बदकिस्मत
\VUgras{मुकाबिला करे वाला} \textsuperscript{(30)} पानी में गिर गइल बा। ऊ
किनारे आवे खातिर तइरत बा।

%%K t03-c2-10-1 p
१०. --- ओकरा से बड़का जवान लोग \VUgras{नाव के भिड़ंत}
\textsuperscript{(32)} करत बाड़ें, आ ई सब \VUgras{घाट}
\textsuperscript{(33)} लगे होत बा। एह सब खेल के
देखे खातिर बड़हन \VUgras{जनता} \textsuperscript{(31)} \VUgras{पुल}
\textsuperscript{(27)} के \VUgras{जँगला} पर टेकल आ किनारा पर जुटल बा।
कुछ देखे वाला लोग तबो मुड़ के \VUgras{इनाम वाला खंभा}
\textsuperscript{(45)} के खेल देखत बा भा ओह नगरपालिका के जुलूस के
निहारत बा जवन इनाम बाँटे आवत बा।

%%K t03-c2-11-1 p
११. --- पहिले, इहाँ कुछ \VUgras{छोकड़ा} \textsuperscript{(35)}
\VUgras{बजनियाँ} \textsuperscript{(36)} के आगे-आगे चलत बाड़ें ; ओकरा बाद
\VUgras{नगर-मुखिया} \textsuperscript{(37)}, ओकर सहायक आ कस्बा के
\VUgras{नगर-पंचायत} आवत बाड़ें। आखिर में \VUgras{दमकल वाला}
\textsuperscript{(38)} लोग चलत बाड़ें।

%%K t03-tit-12 sub
\VUsaut{5.0mm}
\VUpk{10.2pt}{40}{मेला।}
\VUsaut{3.6mm}

%%K t03-c2-12-1 p
१२. --- \VUgras{मेला के मैदान} \textsuperscript{(39)} पर बहुते भीड़ बा।
लोग तरह-तरह के \VUgras{तंबू} देखे आवेला : \VUgras{लकड़ी के घोड़ा}
\textsuperscript{(40)}, \VUgras{सर्कस} \textsuperscript{(41)} जहाँ खेल
शुरू हो चुकल बा, आ \VUgras{सरबजनिक नाच} \textsuperscript{(42)} जहाँ
सहर के जवान लइका आ लइकी लोग \VUgras{नाचे} के मजा लेत बाड़ें।

%%K t03-c2-13-1 p
१३. --- छोर पर हमनी के \VUgras{अखाड़ा} \textsuperscript{(44)} देखत बानी
जा जहाँ बहादुर स्पेनी \VUgras{मातादोर} गुस्साइल \VUgras{साँड़}
\textsuperscript{(45)} से भिड़त बा, ओकरा बाद \VUgras{रेल-झूला}
\textsuperscript{(49)}, आ \VUgras{निसानाबाजी के जगहा} \textsuperscript{(46)}
जहाँ लोग \VUgras{बंदूक} चलावे आ \VUgras{निसाना} लगावे के अभ्यास करेला।

%%K t03-c2-14-1 p
१४. --- ओकरा लगे, इहाँ \VUgras{मेला के नाटकघर} \textsuperscript{(47)}
बा जहाँ \VUgras{हास-नाटक} आ \VUgras{दुखांत नाटक} खेलल जाला। ओकरे लगे
\VUgras{दैत्य} \textsuperscript{(48)} आ \VUgras{बउना} के तंबू बा।
पहिलका के \VUgras{कद} दू मीटर पंद्रह सेंटीमीटर बा ; दुसरका मुसकिल से
एगो बच्चा जेतना बा।

%%K t03-c2-15-1 p
१५. --- आखिर में, इहाँ बड़का \VUgras{जनावरखाना} \textsuperscript{(50)}
बा आपन \VUgras{जंगली जनावर} के साथे — आपन \VUgras{बाघ}
\textsuperscript{(51)} तेज \VUgras{नाखून} के साथे, आपन \VUgras{सिंह}
\textsuperscript{(52)} घन \VUgras{गरदन के बाल} के साथे, आपन दू गो
\VUgras{जिराफ} \textsuperscript{(53)} आ आपन \VUgras{हाथी}
\textsuperscript{(54)} जे आपन \VUgras{सूँड़} एतना चतुराई से चलावेला।

%%K t03-c2-16-1 p
१६. --- ओह जनावरन के सधावे वाला \VUgras{जनावर सधावे वाला}
\textsuperscript{(55)} तंबू के केवाड़ लगे बा।
